\originalchapter{色散、波包与传播速度}\label{ch:disp}
\originalsection{单色相位与缓变包络}
理想柔软弦上的脉冲形状不变, 钢琴弦的刚性却使短波长成分较快, 混合脉冲的不同成分会逐渐错开. 单色波某个相位点满足 $kx-\omega t=\text{常数}$, 其移动速度是相速度 $v_p=\omega/k$. 波包包络的速度通常不同.

把两个相近波数的波相加:
\[
\cos(k_1x-\omega_1t)+\cos(k_2x-\omega_2t)
=2\cos\left(\frac{\Delta k x-\Delta\omega t}{2}\right)
\cos(\bar kx-\bar\omega t).
\]
慢变化余弦是包络, 快变化余弦是载波. 令波数间隔趋小, 包络速度 $v_g=\dd\omega/\dd k$, 称为群速度. 这里 $\Delta\omega$ 与 $\Delta k$ 必须用同一顺序相减. 一般窄带波包写为 $\int A(k)\e^{\ii[kx-\omega(k)t]}\dd k$, 在 $k_0$ 附近展开
\[
\omega(k)=\omega_0+v_g(k-k_0)+\tfrac12\omega''(k_0)(k-k_0)^2+\cdots.
\]
保留一阶时, 积分化为 $\e^{\ii(k_0x-\omega_0t)}F(x-v_gt)$. 不变形近似还要求 $|\omega''(k_0)|t(\Delta k)^2\ll1$, 并检查更高阶项. 窄带只是有利条件, 传播足够远仍会累积展宽. 有强吸收或反常色散时, 波峰群速度不能自动等同于能量或信号前沿速度.

\originalsection{有刚性的弦}
若弦的弯曲刚度引入正系数 $\alpha$（量纲为长度平方）, 模型为 $y_{tt}=c^2(y_{xx}-\alpha y_{xxxx})$. 平面波代入后得
\[
\omega^2=c^2k^2(1+\alpha k^2),\quad
v_p=c\sqrt{1+\alpha k^2},\quad
v_g=c\frac{1+2\alpha k^2}{\sqrt{1+\alpha k^2}}\quad(k>0).
\]
短波受到较大弯曲回复力, 所以频率高于理想弦. 第四阶空间方程需要每端两个边界条件. 简支端固定位置但不传递弯矩, 因弯矩正比于曲率而可取 $y=0,y_{xx}=0$, 此时 $\sin(n\pi x/L)$ 仍是模, $\omega_n=cn\pi\sqrt{1+\alpha n^2\pi^2/L^2}/L$. 若为夹持端, 位置与小转角均被固定, 则 $y=0,y_x=0$, 模形状与允许波数也改变. 因此“边界只决定 $k$、色散只决定 $\omega$”需限定在共享同一空间模及正确边界的模型中.

无色散两端固定弦的所有频率是基频整数倍, 经 $2L/c$ 全部模相位同时回到起点. 有色散时频率比一般非有理, 不存在精确的有限共同周期. 对只有几个频率的运动, 只有它们可公度时才有精确重复; 不能对任意实周期谈通常整数意义的最小公倍数.

\originalsection{截止与不传播的频带}
设某理想介质的色散为 $\omega^2=\omega_b^2+c_m^2k^2$, 其中 $\omega_b$ 是截止角频率. 实波数支上 $v_p=\omega/k=c_m\sqrt{1+\omega_b^2/(c_m^2k^2)}$, $v_g=c_m^2k/\omega$, 所以 $v_pv_g=c_m^2$. 当 $\omega\downarrow\omega_b$ 时 $v_g\to0$, $v_p\to\infty$, 并不表示能量无限快传输. $\omega<\omega_b$ 时 $k=\ii\kappa$, $\kappa=\sqrt{\omega_b^2-\omega^2}/c_m$, 波在介质内指数衰减, 称为倏逝场. 无损半无限介质中这类场没有向深处的平均能量通量, 可导致全反射; 有限厚度层仍可能经另一界面耦合产生透射.

对透明介质的 $k=n(\omega)\omega/c$, 可直接求 $v_g=c/[n+\omega\,\dd n/\dd\omega]$, $v_p=c/n$. 折射率的值与其频率导数分别控制这两种速度. 给出图像色散曲线时, $\omega/k$ 是原点到曲线点的割线斜率, $\dd\omega/\dd k$ 是该点切线斜率.
